\subsection{Hilbertian direct sum and tensor product of unitary representations}

Let G be a topological group. Let \((\varrho_{i})_{i\in\mathrm{I}}\) be a family of unitary representations of G in Hilbert spaces \(E_{i}\) . Let E be the space of the external Hilbert sum of the \(E_{i}\) . For every g in G and for every \(x=(x_{i})_{i\in\mathrm{I}}\) in E, we have

\[
\sum_ {i} \| \varrho_ {i} (g) x _ {i} \| ^ {2} = \sum_ {i} \| x _ {i} \| ^ {2} = \| x \| ^ {2},
\]

which proves that \((\varrho_i(g)x_i)_{i\in \mathrm{I}}\) is in E and has the same norm as \(x\) . Thus the element \(\varrho (g)\colon (x_{i})_{i\in \mathrm{I}}\mapsto (\varrho_{i}(g)x_{i})_{i\in \mathrm{I}}\) is a unitary element of \(\mathcal{L}(\mathrm{E})\) .

Lemma 7. --- The map \(g \mapsto \varrho(g)\) is a unitary representation of G in E.

By lemma 4 of V, p. 380, it suffices to show that for every i in I and every x in \(E_{i}\) , the map \(g \mapsto \varrho_{i}(g)x\) is continuous at the identity element e of G. This map is the composition of the continuous map \(g \mapsto \varrho_{i}(g)x\) and of the canonical injection of \(E_{i}\) into E. It is therefore continuous.

The representation \(\varrho\) is called the Hilbert sum of the unitary representations \((\varrho_{i})_{i\in I}\) ; it is denoted by \(\varrho = \bigoplus \varrho_{i}\) .

Let G and H be topological groups. Let \(\varrho_{1}\) (resp. \(\varrho_{2}\) ) be a unitary representation of G (resp. H) in a Hilbert space E \(_{1}\) (resp. E \(_{2}\) ). Let E = E \(_{1}\) \(\widehat{\otimes}_{2}\) E \(_{2}\) be the Hilbert tensor product space of E \(_{1}\) and E \(_{2}\) (EVT, V, p. 28, def. 1). For \((g,h)\in\mathrm{G}\times\mathrm{H}\) , let \(\varrho(g,h)\) be the continuous endomorphism \(\varrho_{1}(g)\widehat{\otimes}_{2}\varrho_{2}(h)\) of E (TVS, V, p. 28); we will denote it simply \(\varrho_{1}(g)\otimes\varrho_{2}(h)\) whenever no ambiguity is to be feared.

Lemma 8. --- a) The map \(\varrho\colon(g,h)\mapsto\varrho(g,h)\) is a unitary representation of G × H in E;

\begin{enumerate}
\def\labelenumi{\alph{enumi})}
\setcounter{enumi}{1}
\item
  For any orthonormal basis \((e_i)_{i\in \mathrm{I}}\) of \(\mathrm{E}_1\) the mapping of the Hilbert sum \(\bigoplus_{i\in \mathrm{I}}\mathrm{E}_2\) to E defined by \((y_i)_{i\in \mathrm{I}}\mapsto \sum_{i\in \mathrm{I}}e_i\otimes y_i\) is an isometric H-isomorphism of the Hilbert sum \(\bigoplus_{i\in \mathrm{I}}\varrho_2\) in \(\operatorname {Res}_{\{e\} \times \mathrm{H}}^{\mathrm{G}\times \mathrm{H}}(\varrho)\) ;
\item
  For any orthonormal basis \((f_j)_{j\in \mathrm{J}}\) of \(\mathrm{E}_2\) the mapping of the Hilbert sum \(\bigoplus_{j\in \mathrm{J}}\mathrm{E}_1\) to E defined by \((x_{j})_{j\in \mathrm{J}}\mapsto \sum_{j\in \mathrm{J}}x_{j}\otimes f_{j}\) is an isometric G-isomorphism of the Hilbert sum \(\bigoplus_{j\in \mathrm{J}}\varrho_1\) in \(\operatorname {Res}_{\mathbf{G}\times \{e\}}^{\mathbf{G}\times \mathbf{H}}(\varrho)\) .
\end{enumerate}

The map \((g,h)\mapsto\varrho(g,h)\) is a homomorphism from \(\mathrm{G}\times\mathrm{H}\) to \(\mathbf{GL}(\mathrm{E})\) (cf. EVT, V, p. 28, no. 2). Let \((e_{i})_{i\in\mathrm{I}}\) be an orthonormal basis of \(\mathrm{E}_{1}\) . By EVT, V, p. 29, prop. 3 and cor. 2, the map

\[
u \colon (y _ {i}) \mapsto \sum_ {i \in I} e _ {i} \otimes y _ {i}
\]

is an isometry from the Hilbert sum \(F_{2} = \bigoplus_{i \in I} E_{2}\) onto E. Via this isometry, the representation \(\varrho_{\mathrm{H}} \colon h \mapsto \varrho(e, h)\) of H in E is identified with the direct sum representation of the representations \((\varrho_{2})_{i\in I}\) in \(F_{2}\) . In particular, it is a unitary representation of H in E (lemma 7). Likewise, the homomorphism \(\varrho_{G}: g \mapsto \varrho(g,e)\) is a unitary representation of G in E.

Let \((g,h)\in\mathrm{G}\times\mathrm{H}\) . We have \(\varrho(g,h)=\varrho_{\mathrm{G}}(g)\circ\varrho_{\mathrm{H}}(h)\) , hence \(\varrho(g,h)\) is unitary; moreover, \(\varrho_{\mathrm{G}}(g)\) and \(\varrho_{\mathrm{H}}(h)\) commute, so lemma 1 of V, p. 377 implies that \(\varrho\) is a unitary representation of \(G\times H\) .

Finally, the assertions concerning the restriction of \(\varrho\) to the subgroups \(\{e\} \times \mathrm{H}\) and \(\mathrm{G} \times \{e\}\) were obtained in the course of the preceding argument.

The representation \((g,h)\mapsto\varrho_{1}(g)\otimes\varrho_{2}(h)\) of G×H is called the external tensor product of the unitary representations \(\varrho_{1}\) and \(\varrho_{2}\) , and is denoted \(\varrho_{1}\boxtimes\varrho_{2}\) .

Let \(n \in \mathbf{N}\) and let \((\mathrm{G}_i)_{1 \leqslant i \leqslant n}\) be a finite family of topological groups. Let \(\varrho_i\) be a unitary representation of \(\mathrm{G}_i\) in a Hilbert space \(\mathrm{E}_i\) for \(1 \leqslant i \leqslant n\) . One defines similarly a representation

\[
\varrho_ {1} \boxtimes \dots \boxtimes \varrho_ {n}
\]

of \(G_{1} \times \cdots \times G_{n}\) in the Hilbert space \(E = E_{1} \widehat{\otimes}_{2} \cdots \widehat{\otimes}_{2} E_{n}\) (EVT, V, p. 27).

Suppose that \(G_{i} = G\) for \(1 \leqslant i \leqslant n\) . Let \(\Delta_{n}: G \to G^{n}\) be the homomorphism defined by \(g \mapsto (g, \ldots, g)\) for all \(g \in G\) . We denote by \(\varrho_{1} \otimes \cdots \otimes \varrho_{n}\) the unitary representation \((\varrho_{1} \boxtimes \cdots \boxtimes \varrho_{n}) \circ \Delta_{n}\) of G. This is called the tensor product of the unitary representations \(\varrho_{i}\) .

For every permutation \(\sigma\) of \(\{1,\ldots,n\}\) , the canonical isomorphism

\[
\mathrm{E} _ {1} \widehat {\otimes} _ {2} \dots \widehat {\otimes} _ {2} \mathrm{E} _ {n} \rightarrow \mathrm{E} _ {\sigma (1)} \widehat {\otimes} _ {2} \dots \widehat {\otimes} _ {2} \mathrm{E} _ {\sigma (n)}
\]

(EVT, V, p. 28) is an isometric isomorphism

\[
\varrho_ {1} \otimes \dots \otimes \varrho_ {n} \rightarrow \varrho_ {\sigma (1)} \otimes \dots \otimes \varrho_ {\sigma (n)}
\]

of representations of G.

If \(\varrho_{i}=\varrho\) for \(1\leqslant i\leqslant n\) , where \(\varrho\) is a unitary representation of G, we also denote by \(\varrho^{\otimes n}\) the tensor product representation of the \(\varrho_{i}\) , and we say that it is the \(n\)th tensor power of \(\varrho\) .

